\documentclass[twocolumn]{revtex4-2}

\usepackage{amsmath,amssymb}
\usepackage{graphicx}
\usepackage{booktabs}
\usepackage{xcolor}
\usepackage{hyperref}

\newcommand{\CFL}{\textnormal{CFL}}
\newcommand{\QCD}{\textnormal{QCD}}
\newcommand{\Zthree}{\mathbb{Z}_3}

% Auto-generated by update_counts.py — 2026-04-24T19:44:03.140400
% Do not edit manually. Run: uv run python scripts/update_counts.py
\newcommand{\totaltheorems}{3993}
\newcommand{\substantivetheorems}{3882}
\newcommand{\placeholdertheorems}{111}
\newcommand{\axiomcount}{1}
\newcommand{\sorrycount}{0}
\newcommand{\sorrypercent}{0.0}
\newcommand{\leanmodules}{167}
\newcommand{\leandefinitions}{3297}
\newcommand{\pythonmodules}{68}
\newcommand{\testfiles}{59}
\newcommand{\totaltests}{2836}
\newcommand{\figurecount}{111}
\newcommand{\notebookcount}{54}
\newcommand{\papercount}{19}
\newcommand{\aristotleproved}{322}
\newcommand{\aristotleruns}{44}
% --- Paper 18 / Phase 5t per-section counts (FermiHubbardDimer.lean) ---
\newcommand{\fhdTotal}{143}
\newcommand{\fhdWaveTwo}{13}
\newcommand{\fhdWaveThree}{6}
\newcommand{\fhdWaveFour}{22}
\newcommand{\fhdWaveFive}{22}
\newcommand{\fhdWaveSixABC}{38}
\newcommand{\fhdWaveSixExtra}{15}
\newcommand{\fhdWaveSeven}{19}
\newcommand{\fhdWaveEight}{8}


\begin{document}

\title{CFL chiral Lagrangian, emergent $\Zthree$ one-form symmetry, and
the Phase~6d correctness-push: $\CFL$ emergent $\Zthree$ $\equiv$
$\QCD$ center $\Zthree$ as cross-derivation consistency check}

\author{John G. Roehm}
\affiliation{SK-EFT Hawking Project (Phase 6d Wave 3)}

\date{\today}

\begin{abstract}
We formalize Son's color-flavor-locked (CFL) chiral Lagrangian and
the Hirono--Tanizaki emergent $\Zthree$ one-form symmetry of the CFL
phase in Lean 4.29. The Phase~6d.3 \emph{correctness-push} anchor
ships as the load-bearing theorem
\verb|CFL_emergent_Z3_matches_QCD_center_Z3|: the CFL emergent
$\Zthree$ generator (defined independently as
$\omega = e^{2\pi i / 3}$) coincides with the $\QCD$ center $\Zthree$
generator from Phase~6d Wave~1 (\verb|CenterSymmetryConfinement.Z3|).
Two structurally independent derivations --- bare-gauge center
symmetry of UV $\QCD$ vs.\ emergent diquark-sector one-form symmetry
of the CFL IR phase --- yield the \emph{same} generator. Consequences
$\omega^3 = 1$ and $|\omega| = 1$ derive THROUGH the cross-module
identification (load-bearing calls to W1's \verb|centerPhase_pow_N|
and \verb|centerPhase_norm_one|), not via direct $\verb|Complex.exp_nat_mul|$.
The cube-root sum $1 + \omega + \omega^2 = 0$ is the algebraic
fingerprint distinguishing $\Zthree$ from $\Zthree[2]$.
The CFL chiral Lagrangian skeleton (kinetic + mass terms) is encoded
with the chiral-limit Goldstone behavior ($m_q = 0$ vanishes
the mass term) and the in-CFL-phase positivity (requires BOTH
$m_q > 0$ AND $\verb|isCFLPhase|$). The Hirono--Tanizaki
topological-order-beyond-Landau-Ginzburg framing is encoded as a
tracked-hypothesis predicate over $\Zthree$ charges $\{0, 1, 2\}$
with an existence witness (charge $= 1$) and two falsifiers
(trivial $0$ and out-of-range $3$). A substantive cross-bridge to
Wave~2 (\verb|cfl_phase_with_gmor_dual_broken|) consumes BOTH the
W2 \verb|chiral_unbroken_violates_gmor| contrapositive AND the
W3-internal \verb|isCFLPhase_iff_magnitude_pos|, encoding the
\emph{simultaneous} chiral-SSB and CFL-phase breaking. The module
ships $13$ substantive Lean theorems with $0$ sorry and $0$ new
axioms, a Python mirror with $17$ pytest cross-checks ($17/17$ pass)
including a numerical check that the two $\Zthree$ generators agree
to machine precision via independent code paths.
\end{abstract}

\maketitle

\section{Motivation: the cross-derivation consistency check}
\label{sec:motivation}

The Phase~6d roadmap identifies $\CFL$ as the correctness-push anchor
of the Phase~6d QCD-interpretation layer~\cite{Phase6dRoadmap}. The
specific construct is the Hirono--Tanizaki framing~\cite{HironoTanizaki2018}:
the CFL phase exhibits an emergent $\Zthree$ one-form symmetry that
should match the $\QCD$ center $\Zthree$ from Wave~1
(\verb|CenterSymmetryConfinement|). \emph{Either outcome is
substantive}: agreement provides an independent consistency check on
the project's emergent-symmetry identification; disagreement would
imply structural tension.

This wave ships the agreement branch, formalized as a direct
identification at the level of the cyclic generator:
\begin{equation}
  \omega_\CFL := e^{2\pi i / 3} = \verb|centerPhase|(\verb|Z3|).
  \label{eq:correctness-push}
\end{equation}

The substantive content is not algebraic novelty (both quantities
compute trivially to $\cos(2\pi/3) + i \sin(2\pi/3)$); it is the
cross-module identification across two derivation chains:
\begin{enumerate}
\item bare-gauge center symmetry of $\QCD$ via the
  \verb|CenterZN| structure (W1, $N = 3$);
\item emergent one-form symmetry of the CFL diquark sector via the
  Hirono--Tanizaki framing (this wave).
\end{enumerate}

\section{CFL diquark order parameter}
\label{sec:diquark}

The CFL phase is encoded as a complex diquark VEV
\begin{verbatim}
abbrev CFLDiquark : Type := C
def isCFLPhase (Phi : CFLDiquark) : Prop :=
  Phi != 0
\end{verbatim}
with the order parameter
\begin{equation}
  \verb|diquarkMagnitude|(\Phi) := \|\Phi\|.
\end{equation}
The qualitative-quantitative bridge
\verb|isCFLPhase_iff_magnitude_pos| ($\verb|isCFLPhase|\,\Phi
\Leftrightarrow 0 < \|\Phi\|$) parallels W1's
\verb|confining_iff_magnitude_zero| pattern, applied dually
(``$|\Phi| > 0$ in the CFL phase'' vs.\ ``$|P| = 0$ in the
confining phase'').

\section{Emergent $\Zthree$ one-form symmetry}
\label{sec:z3}

\subsection{Independent definition}

The emergent $\Zthree$ generator is defined \emph{independently} from
W1's $\verb|centerPhase|(\verb|Z3|)$:
\begin{verbatim}
noncomputable def emergentZ3Phase : C :=
  cexp (2 * pi * I / 3)
\end{verbatim}
This independence is critical: the definition does not refer to W1
via aliasing; the equality is proven, not assumed.

\subsection{The correctness-push theorem}

\begin{verbatim}
theorem CFL_emergent_Z3_matches_QCD_center_Z3 :
    emergentZ3Phase =
    SKEFTHawking.CenterSymmetryConfinement.centerPhase
      SKEFTHawking.CenterSymmetryConfinement.Z3 := by
  unfold emergentZ3Phase
  unfold SKEFTHawking.CenterSymmetryConfinement.centerPhase
  norm_num [SKEFTHawking.CenterSymmetryConfinement.Z3]
\end{verbatim}
The proof is short (both definitions reduce to $\verb|cexp|(2\pi i / 3)$
after $\verb|Z3|.N = 3$ unfolds), but the substantive content is the
\emph{identification across two independent definitions}.

\subsection{Consequences via cross-module calls}

The cube-root and unit-modulus consequences derive THROUGH the
identification, using W1's load-bearing theorems:
\begin{align}
  \omega_\CFL^3 &= 1 &&\text{(via \texttt{centerPhase\_pow\_N})} \\
  \|\omega_\CFL\| &= 1 &&\text{(via \texttt{centerPhase\_norm\_one})}
\end{align}

The cube-roots-sum identity is the algebraic distinguishing feature:
\begin{equation}
  1 + \omega_\CFL + \omega_\CFL^2 = 0
  \label{eq:cube-sum}
\end{equation}
with the proof using the factorization $X^3 - 1 = (X - 1)(X^2 + X + 1)$
at $X = \omega$, plus $\omega \neq 1$ derived from the imaginary part
$\sin(2\pi/3) > 0$.

Figure~\ref{fig:cfl} (left panel) visualizes the three cube roots
$\{1, \omega, \omega^2\}$ on the unit circle.

\begin{figure*}[ht]
  \centering
  \includegraphics[width=0.92\linewidth]{figures/fig_cfl_z3_center_bridge.png}
  \caption{
    \emph{Phase 6d Wave 3 — CFL emergent $\Zthree$ (Hirono--Tanizaki)
    $\equiv$ QCD center $\Zthree$ (W1) correctness-push
    identification.} Left: cube roots of unity $\{1, \omega,
    \omega^2\}$ on the complex unit circle, with the CFL emergent
    generator (independent definition in this wave) and the QCD
    center generator from \texttt{CenterSymmetryConfinement.Z3} BOTH
    at $\omega = e^{2\pi i / 3}$. The sum rule $1 + \omega + \omega^2
    = 0$ is the algebraic fingerprint distinguishing $\Zthree$ from
    $\Zthree[2]$. Right: bar chart classifying the four candidate
    charges $\{0, 1, 2, 3\}$ by the Hirono--Tanizaki topological-
    order-beyond-Landau-Ginzburg sector. Charges $1$ and $2$ are
    topological witnesses; charges $0$ and $3$ are falsifiers.
  }
  \label{fig:cfl}
\end{figure*}

\section{CFL chiral Lagrangian skeleton}
\label{sec:lagrangian}

The Lagrangian sector ships at the structural level only — the
canonical kinetic term and the leading mass term:
\begin{align}
  \verb|cflKineticTerm|(\|\partial \Phi\|) &:= \tfrac{1}{2}
    \|\partial \Phi\|^2, \\
  \verb|cflMassTerm|(m_q, \Phi) &:= m_q \cdot \|\Phi\|^2.
\end{align}
The kinetic term non-negativity follows by $\verb|positivity|$. The
mass term has two key properties:
\begin{itemize}
\item \verb|cflMassTerm_chiral_limit|: vanishes at $m_q = 0$ (the
  Goldstone-boson limit of CFL pseudoscalars).
\item \verb|cflMassTerm_pos_in_cfl_phase|: strictly positive when BOTH
  $m_q > 0$ AND $\verb|isCFLPhase|\,\Phi$ — the proof uses the
  qualitative-quantitative bridge to convert $\Phi \neq 0$ to
  $\|\Phi\| > 0$.
\end{itemize}

\section{Topological order beyond Landau-Ginzburg}
\label{sec:topological}

The Hirono--Tanizaki framing is encoded as a tracked-hypothesis
predicate over $\Zthree$ charges:
\begin{verbatim}
def H_TopologicalOrderBeyondLG (charge : N) :
  Prop :=
    charge < 3 AND charge != 0
\end{verbatim}
with two independent constraints (drop-conjunct test passes):
cyclicity ($\charge < 3$, restricting to $\Zthree$) and
non-triviality ($\charge \neq 0$, excluding the vacuum sector).

The witness theorem \verb|H_TopologicalOrderBeyondLG_witness| ships
$\charge = 1$ as a satisfying instance; two falsifier theorems show
that the trivial vacuum ($\charge = 0$) and out-of-range
($\charge = 3$) sectors fail the predicate.

\section{Cross-bridge to Wave 2}
\label{sec:bridges}

The W2 cross-bridge consumes both upstream and W3-internal load-bearing
content:
\begin{verbatim}
theorem cfl_phase_with_gmor_dual_broken
    (Phi : CFLDiquark) (h_cfl : isCFLPhase Phi)
    (m_pi f_pi m_q sigma : R)
    (h_mq : 0 < m_q) (h_pi : m_pi != 0)
    (h_fpi : f_pi != 0)
    (h_gmor : m_pi^2 * f_pi^2 = -2 * m_q * sigma) :
    sigma < 0 AND 0 < diquarkMagnitude Phi
\end{verbatim}
The proof uses W2's \verb|chiral_unbroken_violates_gmor| (for the
$\sigma < 0$ conjunct) AND W3's \verb|isCFLPhase_iff_magnitude_pos|
(for the $\|\Phi\| > 0$ conjunct). Both conjuncts are independently
load-bearing. The original first-pass theorem
\verb|cfl_phase_implies_chiral_broken_via_gmor| had \verb|h_cfl| as
an unused parameter (decorative); Lean's unused-variable linter caught
this, and the strengthened form uses the CFL hypothesis substantively.

\section{Discipline metric and Phase 6d closure}
\label{sec:discipline}

The preemptive-strengthening discipline + multi-pass-review protocol
(adopted from 6d.1 finding) was applied. First-pass writing caught
ONE pattern (the unused-\verb|h_cfl| identity-wrapper, fixed in
session). Subsequent pass found nothing. Multi-pass-protocol
terminates: $1$ retroactive theorem.

\paragraph*{Phase 6d trend.}
\begin{tabular}{l|l|l}
Wave & Total retroactive & First-pass cost\\\hline
6c.3 & 12 & 12 (baseline)\\
6b.1 & 5 & 5\\
6d.1 & 6 & 3\\
6d.2 & 4 & 1\\
6d.3 (this) & 1 & 1\\
\end{tabular}

\paragraph*{Phase 6d closure.}
With this wave, all three Track waves ship: W1
(\verb|CenterSymmetryConfinement|), W2 (\verb|ChiralSSB_QCD|), W3
(\verb|CFLChiralLagrangian|). The roadmap correctness-push (CFL
emergent $\Zthree$ $\equiv$ QCD center $\Zthree$) ships in the
agreement branch. Independent next steps within Phase~6d are absorbed
by HepLean / PhysLean per the roadmap scope lock.

\begin{thebibliography}{99}

\bibitem{AlfordRajagopalWilczek1999}
M.~G.~Alford, K.~Rajagopal, and F.~Wilczek,
\emph{Color-flavor locking and chiral symmetry breaking in high
density QCD}, Nucl. Phys. B \textbf{537}, 443 (1999).

\bibitem{SonStephanov2001}
D.~T.~Son and M.~A.~Stephanov,
\emph{Inverse meson mass ordering in the color-flavor-locking phase
of high-density QCD}, Phys. Rev. Lett. \textbf{86}, 592 (2001).

\bibitem{SchaeferWilczek1999}
T.~Sch\"afer and F.~Wilczek,
\emph{Continuity of quark and hadron matter}, Phys. Rev. Lett.
\textbf{82}, 3956 (1999).

\bibitem{HironoTanizaki2018}
Y.~Hirono and Y.~Tanizaki,
\emph{Quark-hadron continuity beyond the Landau-Ginzburg paradigm},
JHEP \textbf{12}, 062 (2018).

\bibitem{Phase6dRoadmap}
\emph{Phase 6d Roadmap} \verb|docs/roadmaps/Phase6d_Roadmap.md|
(SK-EFT Hawking project, 2026).

\bibitem{CenterSymmetryConfinementLean}
\emph{Lean module} \verb|lean/SKEFTHawking/CenterSymmetryConfinement.lean|
(SK-EFT Hawking project, Phase 6d Wave 1, 2026).

\bibitem{ChiralSSBQCDLean}
\emph{Lean module} \verb|lean/SKEFTHawking/ChiralSSB_QCD.lean|
(SK-EFT Hawking project, Phase 6d Wave 2, 2026).

\bibitem{Hatsuda2008}
T.~Hatsuda, M.~Tachibana, N.~Yamamoto, and G.~Baym,
\emph{New critical point induced by the axial anomaly in dense QCD},
Phys. Rev. Lett. \textbf{97}, 122001 (2006).

\bibitem{GaiottoKapustinSeibergWillett2015}
D.~Gaiotto, A.~Kapustin, N.~Seiberg, and B.~Willett,
\emph{Generalized global symmetries},
JHEP \textbf{02}, 172 (2015).

\bibitem{Polyakov1978}
A.~M.~Polyakov,
\emph{Thermal properties of gauge fields and quark liberation},
Phys. Lett. B \textbf{72}, 477 (1978).

\end{thebibliography}

\end{document}
